Die Position (9) bedarf ebenfalls noch einiger Erläuterungen. In der Jahresrechnung für 1850—1851 ist das Geldvermögen der Kasse zum 1. Juli 1851 mit 125076 Thl. 1 Ggr. 11 Pf. angesetzt; davon entfallen auf die verzinslichen Kapitalien 123372 Thl. 2 Ggr. 7 Pf. Der Rest besteht aus dem Bargeldbestand, den Rückständen und einem dem Universitätsapotheker bewilligten zinslosen Vorschuss. Die Kapitalien sind etwa zur Hälfte als Hypothekendarlehen an private Schuldner, im Übrigen in unkündbaren Staatspapieren angelegt. Diese sind in der neuesten Jahresrechnung, ebenso wie schon in mehreren früheren, ohne Weiteres zum Nennwert angesetzt. Im Jahr 1845 waren diese Ansätze dagegen nach den Ankaufspreisen gemacht worden; ich habe sie in meiner damaligen Rechnung unverändert beibehalten, weil die Wertschwankungen der Staatspapiere damals viel geringer waren als in den letzten 3 bis 4 Jahren. Jetzt, da ein beträchtlicher Teil der zum Vermögen der Witwenkasse gehörenden Staatspapiere so weit unter dem Nennwert steht, halte ich es für notwendig, die Papiere in der Bilanzrechnung zu dem derzeitigen tatsächlichen Wert anzusetzen, zu dem sie sich verkaufen lassen. Dazu habe ich die Börsenkurse in Frankfurt und Hannover vom 1. Oktober verwendet, weil ein bestimmtes Datum gewählt werden musste. Es sind die folgenden
\begin{center}
\footnotesize\setlength{\tabcolsep}{2pt}
\begin{tabular}{l r|l r}
Österreichische \(4\tfrac{1}{2}\) Metalliques . .&\(67\tfrac{1}{8}\)&Hannoversche 5 Proz. . . .&\(104\tfrac{1}{2}\)\\
— \(4\tfrac{1}{2}\) bei Goll . . . .&\(69\tfrac{3}{4}\)&— 4 Proz. . . .&102\\
— \(4\tfrac{1}{2}\) bei Bethmann&\(71\tfrac{3}{4}\)&— \(3\tfrac{1}{2}\) Courant&\(98\tfrac{1}{4}\)\\
Badische \(3\tfrac{1}{2}\) . . . . . . . . . . . . . .&\(87\tfrac{3}{4}\)&— \(3\tfrac{1}{2}\) Gold . .&\(97\tfrac{3}{4}\)
\end{tabular}
\end{center}
\SourceLeaf{173}\noindent\hspace*{\parindent}Den Goldkurs habe ich mit 1 Louisd’or \(=\) 9 fl. 38 kr. und \(=5\tfrac{1}{2}\) Thaler Courant angenommen. Einige Kurse sind seitdem übrigens noch etwas, wenn auch nur wenig, gefallen.

Das Ergebnis dieser Umrechnungen ist, dass dieselben verzinslichen Kapitalien, die nach dem Nennwert mit 123372 Thaler angesetzt sind, nach dem derzeitigen Börsenkurswert nur 119044 Thl. 23 Ggr. erbringen. Dieser Summe füge ich noch die 800 Thaler hinzu, für die das frühere Müllersche Haus verkauft worden ist. Mit den anderen Posten (Bargeldbestand usw.) ergibt sich so für das Geldvermögen der Witwenkasse zum 1. Juli 1851
\begin{center}121548 Thaler 22 Ggr.\end{center}
\noindent\hspace*{\parindent}Um, soweit ich dazu imstande bin, die Umrechnung auf den 1. Oktober 1851 abzuschätzen, ziehe ich von dieser Summe die zu Michaelis fällig gewordenen Witwenpensionen mit . 1787 Thl. 12 Ggr. ab und füge hinzu
\begin{center}
\small\setlength{\tabcolsep}{2pt}
\begin{tabular}{l r l r l}
\shortstack[l]{den Betrag der Beiträge für das Jahr 1850 bis 1851\\. . . . . . . . . . . . . . . . . . . . . .}&780&,,&&\\
\shortstack[l]{den halbjährlichen Pachtzins für die Apotheke\\. . . . . . . . . . . . . . . . . . . . . . . . . . .}&500&,,&&\\
\shortstack[l]{und einen vierteljährlichen Betrag der Zinseinnahmen,\\auf Grundlage der letztjährigen Einnahmen nach Abzug\\der Einziehungskosten . . . . . . . . . . . . . . . . . . . . . .}&1248&,,&10&,,
\end{tabular}
\end{center}
\noindent woraus sich dann die obige Position (9) ergibt.

\begin{center}\rule{.28\linewidth}{.3pt}\end{center}
